\chapter{1958年沪闽卷}

\section{解答题}

\begin{problem}
设方程 \(x^{2}-(m-2)x-(m-6)=0\) 有二个不相等实根，决定 \(m\) 范围.
\end{problem}

\begin{solution}
该方程为关于 \(x\) 的一元二次方程，其判别式为
\[
\Delta=\left[-(m-2)\right]^{2}-4\times 1\times\left[-(m-6)\right]
=(m-2)^{2}+4(m-6)=m^{2}-20
\]
方程有两个不相等的实根，当且仅当 \(\Delta>0\)，所以 \(m^{2}-20>0\)，即 \(m^{2}>20\). 因而 \(m<-\sqrt{20}\) 或 \(m>\sqrt{20}\). 故 \(m\) 的取值范围为 \(\left(-\infty,-\sqrt{20}\right)\cup\left(\sqrt{20},+\infty\right)\).
\end{solution}

\begin{problem}
解方程组
\(
\begin{cases}
\dfrac{1}{2}\lg x+\dfrac{1}{2}\lg y-\lg\left(4-\sqrt{x}\right)=0,\\
\left(25^{\sqrt{x}}\right)^{\sqrt{y}}-125\times 5^{\sqrt{y}}=0.
\end{cases}
\)
\end{problem}

\begin{solution}
由对数式有意义，必须满足 \(x>0\)，\(y>0\)，且 \(4-\sqrt{x}>0\). 第一个方程化为
\(\lg\frac{\sqrt{x}\cdot\sqrt{y}}{4-\sqrt{x}}=\lg 1\)，
所以
\(\frac{\sqrt{x}\cdot\sqrt{y}}{4-\sqrt{x}}=1\)，
从而
\(\sqrt{x}\left(1+\sqrt{y}\right)=4\)，
即
\(\sqrt{x}=\frac{4}{1+\sqrt{y}}\).

第二个方程中，\(25=5^{2}\)，\(125=5^{3}\)，故
\(5^{2\sqrt{x}\cdot\sqrt{y}}=5^{3+\sqrt{y}}\).
由于底数 \(5>1\)，指数相等，即
\(2\sqrt{x}\cdot\sqrt{y}=3+\sqrt{y}\).
令 \(t=\sqrt{y}\)，则 \(t>0\)，代入上式中的 \(\sqrt{x}=\frac{4}{1+t}\)，得
\(\frac{8t}{1+t}=3+t\).
因 \(1+t>0\)，两边同乘 \(1+t\)，得
\[
8t=(3+t)(1+t)=3+4t+t^{2}
\]
整理得 \(t^{2}-4t+3=0\)，即 \((t-1)(t-3)=0\). 因此 \(t=1\) 或 \(t=3\).

当 \(t=1\) 时，\(y=1\)，\(\sqrt{x}=4/(1+1)=2\)，所以 \(x=4\). 当 \(t=3\) 时，\(y=9\)，\(\sqrt{x}=4/(1+3)=1\)，所以 \(x=1\).

检验：当 \((x,y)=(4,1)\) 时，\(\sqrt{x}\cdot\sqrt{y}=2=4-\sqrt{x}\)，且两个幂的指数均为 \(4\)；当 \((x,y)=(1,9)\) 时，\(\sqrt{x}\cdot\sqrt{y}=3=4-\sqrt{x}\)，且两个幂的指数均为 \(6\). 两组解都满足定义域条件，故原方程组的解为
\(\left(x,y\right)=\left(4,1\right)\) 或 \(\left(x,y\right)=\left(1,9\right)\).
\end{solution}

\begin{problem}
一九五七年我国钢产量约 \(500\text{ 万吨}\)，同年英国钢产量 \(2000\text{ 万吨}\)，根据英国每年平均增长率为 \(4.24\%\)，如果我国准备四年在钢的产量方面赶上英国，问每年平均增产的百分率是多少（精确到 \(0.01\)）.
\end{problem}

\begin{solution}
设我国钢产量的年平均增长率为 \(X\)，则四年后我国钢产量为 \(500\left(1+X\right)^{4}\text{ 万吨}\)，英国钢产量为 \(2000\left(1+0.0424\right)^{4}\text{ 万吨}\). 由四年后两国钢产量相等，得
\(500\left(1+X\right)^{4}=2000\left(1+0.0424\right)^{4}\).
两边同除以 \(500\left(1.0424\right)^{4}\)，得
\(\left(\frac{1+X}{1.0424}\right)^{4}=4\).
增长率为正，故 \(1+X>0\)，取正的四次方根，得
\(\frac{1+X}{1.0424}=\sqrt{2}\)，所以 \(X=1.0424\sqrt{2}-1\).
计算得 \(X\approx0.474176\)，所以年平均增产率约为 \(47.4176\%\). 按题意精确到 \(0.01\%\)，得 \(X\approx47.42\%\).
\end{solution}

\begin{problem}
设有函数 \(y=ax^{2}+bx+c\)，当 \(x=1\) 时，\(y=-1\)；当 \(x=2\) 时，\(y=1\)；又当 \(x=-2\) 时，\(y=17\).

\begin{enumerate}
\item 决定这函数式；
\item 这函数有无极大值或极小值？如果有，其值是什么？
\item 当 \(x\) 是什么值时，这函数值为 \(0\)？
\item 作出函数图像.
\end{enumerate}
\end{problem}

\begin{solution}
\begin{enumerate}
\item 记 \(f(x)=ax^{2}+bx+c\). 由题意，\(a\)，\(b\)，\(c\) 满足
\[
\begin{cases}
a+b+c=-1,\\
4a+2b+c=1,\\
4a-2b+c=17.
\end{cases}
\]
第二个方程减去第三个方程，得 \(4b=-16\)，所以 \(b=-4\). 第二个方程减去第一个方程，得 \(3a+b=2\)，于是
\[
a=\frac{2-b}{3}=\frac{2+4}{3}=2
\]
将 \(a=2\)，\(b=-4\) 代入第一个方程，得 \(c=-1-a-b=-1-2+4=1\). 故函数式为 \(y=2x^{2}-4x+1\).

\item 将函数式配方，得
\[
y=2x^{2}-4x+1=2\left(x^{2}-2x+1\right)-1=2\left(x-1\right)^{2}-1
\]
因为 \(\left(x-1\right)^{2}\geq 0\)，所以 \(y\geq-1\)，且当且仅当 \(x=1\) 时取等号. 又因 \(x\) 的绝对值可以无限增大，\(y=2\left(x-1\right)^{2}-1\) 没有最大值. 因此函数有极小值 \(-1\)，没有极大值；极小值点为 \(\left(1,-1\right)\)，对称轴为 \(x=1\)，抛物线开口向上.

\item 令 \(y=0\)，得 \(2x^{2}-4x+1=0\). 其判别式为 \(\Delta=(-4)^{2}-4\times2\times1=8\)，所以 \(x=\frac{4\pm\sqrt{8}}{4}=1\pm\frac{\sqrt{2}}{2}\).
故当 \(x=1-\frac{\sqrt{2}}{2}\) 或 \(x=1+\frac{\sqrt{2}}{2}\) 时，函数值为 \(0\).

\item 函数图像是开口向上的抛物线 \(y=2\left(x-1\right)^{2}-1\)，顶点为 \(\left(1,-1\right)\)，对称轴为 \(x=1\)，与 \(x\) 轴交于 \(\left(1-\frac{\sqrt{2}}{2},0\right)\) 和 \(\left(1+\frac{\sqrt{2}}{2},0\right)\)，与 \(y\) 轴交于 \(\left(0,1\right)\). 根据这些特征作图如下.
\solutionfigure{\bitmapfigure[max width=0.30\linewidth]{img/1958/hu_min/q4_solution_fig1.png}}
\end{enumerate}
\end{solution}

\begin{problem}
河的对岸有 \(A\)，\(B\) 两点，要在这边测量 \(A\)，\(B\) 间的距离，应该怎样测量？试拟出一种方法来说明之，并写出表达距离的式子（即把能测得的角和能量得的距离的数值表示 \(A\)，\(B\) 间的距离）.
\end{problem}

\begin{solution}
在河的这边选择地势平坦且能同时看见 \(A\)，\(B\) 的两点 \(C\)，\(D\)，把 \(CD\) 测出，记为 \(\alpha\). 线段 \(CD\) 称为基线. 再测量 \(\angle ACD=a\)，\(\angle BCD=\beta\)，\(\angle BDC=\gamma\)，\(\angle ADC=\theta\).
\solutionfigure{\bitmapfigure[width=0.40\linewidth]{img/1958/hu_min/q5_solution_fig1.png}}
在 \(\triangle ACD\) 中，\(\angle CAD=180^{\circ}-(a+\theta)\)，由正弦定理得
\[
\frac{AD}{\sin a}=\frac{CD}{\sin\left(180^{\circ}-(a+\theta)\right)}=\frac{\alpha}{\sin(a+\theta)}
\]
所以 \(AD=\dfrac{\alpha\sin a}{\sin(a+\theta)}\). 同理，在 \(\triangle BCD\) 中，\(\angle CBD=180^{\circ}-(\beta+\gamma)\)，由正弦定理得
\[
\frac{BD}{\sin\beta}=\frac{CD}{\sin\left(180^{\circ}-(\beta+\gamma)\right)}=\frac{\alpha}{\sin(\beta+\gamma)}
\]
所以 \(BD=\dfrac{\alpha\sin\beta}{\sin(\beta+\gamma)}\).

由图中两条视线在 \(D\) 点形成的夹角，\(\angle ADB=\left|\gamma-\theta\right|\)，故在 \(\triangle ADB\) 中由余弦定理得
\[
\begin{aligned}
AB^{2}&=AD^{2}+BD^{2}-2AD\cdot BD\cos\left(\gamma-\theta\right)\\
&=\alpha^{2}\left(\frac{\sin^{2}a}{\sin^{2}(a+\theta)}+\frac{\sin^{2}\beta}{\sin^{2}(\beta+\gamma)}-\frac{2\sin a\sin\beta\cos\left(\gamma-\theta\right)}{\sin(a+\theta)\sin(\beta+\gamma)}\right).
\end{aligned}
\]
因此，测得基线 \(CD=\alpha\) 和上述四个角后，取上式的正平方根，即可得到 \(AB\).
\end{solution}

\begin{problem}
两平行线 \(L_{1}\) 与 \(L_{2}\) 分别和已知圆相切于 \(A\) 与 \(B\) 点，作已知圆的任一切线 \(L_{3}\)，与 \(L_{1}\)，\(L_{2}\) 分别相交于 \(C\) 与 \(D\) 点，证明乘积 \(AC\cdot BD\) 是常量.
\end{problem}

\begin{solution}
设 \(L_{3}\) 与圆 \(O\) 相切于 \(E\)，连接 \(OA\)，\(OB\)，\(OE\)，\(OC\)，\(OD\). 由于半径垂直于切线，有
\[
\angle OAC=\angle OEC=\angle OED=\angle OBD=90^{\circ}
\]
\solutionfigure{\bitmapfigure[width=0.34\linewidth]{img/1958/hu_min/q6_solution_fig1.png}}
所以四边形 \(OACE\) 的一对对角互补，点 \(O\)，\(A\)，\(C\)，\(E\) 共圆. 于是 \(\angle AOE+\angle ACE=180^{\circ}\). 又因 \(L_{1}\parallel L_{2}\)，所以 \(\angle BDE+\angle ACE=180^{\circ}\)，
从而 \(\angle AOE=\angle BDE\).

在直角三角形 \(AOC\) 与 \(EOC\) 中，\(OA=OE\)，且 \(OC\) 是公共斜边，所以两三角形全等，从而 \(\angle AOC=\angle COE\). 同理，在直角三角形 \(BDO\) 与 \(EDO\) 中，\(OB=OE\)，且 \(OD\) 是公共斜边，所以 \(\angle BDO=\angle ODE\). 于是 \(\angle AOE=2\angle COE\)，\(\angle BDE=2\angle ODE\)，结合 \(\angle AOE=\angle BDE\)，得 \(\angle COE=\angle ODE\)，进而 \(\angle AOC=\angle BDO\).

在直角三角形 \(AOC\) 和 \(BDO\) 中，\(\angle OAC=\angle OBD=90^{\circ}\)，且 \(\angle AOC=\angle BDO\)，所以两三角形相似，从而 \(\frac{AC}{OB}=\frac{OA}{BD}\). 交叉相乘，得 \(AC\cdot BD=OA\cdot OB\). 设已知圆半径为 \(r\)，则 \(OA=OB=r\)，故 \(AC\cdot BD=r^{2}\).
右边只由已知圆决定，与切线 \(L_{3}\) 的位置无关，所以 \(AC\cdot BD\) 是常量.
\end{solution}

\begin{problem}
设平面 \(S\) 上有一直角三角形 \(ABC\)，已知一个直角边长为 \(18\text{ 分米}\)，在平面 \(S\) 外取一点 \(P\)，使它与 \(A\)，\(B\)，\(C\) 三点等距离，且设从 \(P\) 点到平面 \(S\) 的距离为 \(40\text{ 分米}\)，求 \(P\) 点到另一直角边的距离.
\end{problem}

\begin{solution}
如图，不妨将已知直角边记为 \(BC=18\text{ 分米}\)，并设 \(\angle ACB=90^{\circ}\). 设 \(D\) 是点 \(P\) 到平面 \(S\) 的垂足，\(E\) 是点 \(P\) 到另一直角边 \(AC\) 的垂足，则 \(PD=40\text{ 分米}\)，且要求 \(PE\).
\solutionfigure{\bitmapfigure[width=0.42\linewidth]{img/1958/hu_min/q7_solution_fig1.png}}
因为 \(PD\perp S\)，在直角三角形 \(PDA\)，\(PDB\)，\(PDC\) 中分别有 \(PA^{2}=PD^{2}+DA^{2}\)，\(PB^{2}=PD^{2}+DB^{2}\)，\(PC^{2}=PD^{2}+DC^{2}\).
又因 \(PA=PB=PC\)，所以 \(DA=DB=DC\). 这说明 \(D\) 是三角形 \(ABC\) 的外心. 直角三角形的外心是斜边 \(AB\) 的中点，故 \(D\) 是 \(AB\) 的中点.

由 \(PE\perp AC\) 以及 \(PA=PC\)，在直角三角形 \(PEA\) 和 \(PEC\) 中，斜边相等且有公共直角边 \(PE\)，所以 \(EA=EC\). 因而 \(E\) 是 \(AC\) 的中点. 在三角形 \(ABC\) 中，连接两个边中点 \(D\)，\(E\)，则
\(DE\parallel BC\)，且 \(DE=\frac{1}{2}BC=9\text{ 分米}\).
又因 \(PD\perp S\)，而 \(DE\) 在平面 \(S\) 内，所以 \(PD\perp DE\). 在直角三角形 \(PDE\) 中，\(PE\) 为斜边，由勾股定理
\[
PE=\sqrt{PD^{2}+DE^{2}}=\sqrt{40^{2}+9^{2}}=\sqrt{1600+81}=\sqrt{1681}=41\text{ 分米}
\]
故点 \(P\) 到另一直角边 \(AC\) 的距离为 \(41\text{ 分米}\).
\end{solution}
